\documentclass[12pt]{article}
% Préambule commun à tous les documents extraits de « La Revue CAPES »
\usepackage[T1]{fontenc}
\usepackage[utf8]{inputenc}
\usepackage{lmodern}
\usepackage{amsmath,amssymb,amsthm,amsfonts,mathrsfs}
\usepackage{setspace}
\usepackage{listings}
\usepackage{pgfplots}
\pgfplotsset{compat=1.18}
\usepackage{longtable}
\usepackage{adjustbox}
\usepackage{lettrine}
\usepackage{tkz-tab}
\usepackage{changepage}
\usepackage{pdflscape}
\usepackage{graphicx}
\usepackage{tikz}
\usepackage{xcolor}
\usepackage[a4paper,top=2cm,bottom=2.5cm,left=2.5cm,right=2cm]{geometry}
\usepackage{fancyhdr}
\usepackage{draftwatermark}
\SetWatermarkText{La RevueCapes}
\SetWatermarkLightness{0.9}
\SetWatermarkScale{2}
\usepackage{hyperref}
\hypersetup{colorlinks=true,linkcolor=blue!50!black,urlcolor=blue!50!black}

\definecolor{revue}{RGB}{20,60,120}

\pagestyle{fancy}
\fancyhf{}
\renewcommand{\headrulewidth}{0.4pt}
\setlength{\headheight}{15pt}

% \entete{Matière}{Titre}{Type de document}
\newcommand{\entete}[3]{%
  \fancyhead[L]{\small\textsc{La Revue CAPES} -- #1}%
  \fancyhead[R]{\small #2}%
  \fancyfoot[C]{\thepage}%
  \thispagestyle{plain}%
  \begin{center}
    {\color{revue}\rule{\textwidth}{1.2pt}}\\[4pt]
    {\small\textsc{Concours directs CAP-CEG / CAPES -- Burkina Faso}}\\[6pt]
    {\LARGE\bfseries\color{revue} #2}\\[6pt]
    {\large\itshape #1 -- #3}\\[2pt]
    {\color{revue}\rule{\textwidth}{1.2pt}}
  \end{center}
  \vspace{0.5cm}%
}

% Encadré pour les remarques ajoutées dans les corrigés
\newcommand{\remarque}[1]{\par\medskip\noindent\fcolorbox{revue}{revue!6}{\parbox{\dimexpr\linewidth-2\fboxsep-2\fboxrule}{\textbf{\color{revue}Remarque.} #1}}\par\medskip}


\begin{document}
\entete{Mathématiques}{Sujet type n°13}{Énoncé}
\begin{spacing}{1.5}

\medskip\textbf{Exercice 1}

\textbf{Partie 1}

Soit $A$ un réel strictement positif. 

Montrer que, pour tout $n$ entier naturel non nul, on peut déterminer une suite unique de $n+1$ 
nombres réels $(t_0, t_1,\cdots, t_{n-1},t_n)$ vérifiant les trois conditions suivantes : 

(i)    $t_0 = 0, t_n = A$ 
(ii)   $t_0 < t_1 < \cdots<t_{n-1} < t_n$
(iii)  $\forall k\in \lbrace 0, ...., n-1\rbrace, t_{k+1}-t_{k}= A/n$

\textbf{Partie 2}

Soit la fonction $f : [0,1]\longrightarrow \mathbb{R}$, continue et strictement positive. 

1) Montrer que, pour tout $n$ entier non nul, il existe dans $[0,1]^{n+1}$ une suite unique de $n+1$ 
nombres réels $(x_0, x_1,\cdots,x_{n-1},  x_n)$ vérifiant les trois conditions suivantes : 

(i)  $x_0 = 0, x_n = 1$

(ii) $x_0 < x_1 < \cdots<x_{n-1} < x_n$

(iii) $\forall k \in \{0, ...., n-1\}, \int_{x_k}^{x_{k+1}}f(t)dt=\frac{1}{n}\int_{0}^{1}f(t)dt$
 
2) Soit $U(n) = \frac{1}{n}\sum_{k=0}^{n}f(x_k)$ 

Déterminer la limite $L$ de $U(n)$ quand $n$ tend vers $+\infty$. 

3) Calculer $L$ dans le cas particulier de $f(x) = e^x$

\medskip\textbf{Exercice 2}

Le symbole $ln$ représente le logarithme népérien et $\mathbb{R}$ désigne l'ensemble des nombres réels. 

1) Soit l'application $g : ]0,1[\cup ]1,+\infty[\longrightarrow \mathbb{R}$, définie par : 
$g(x)=\frac{1}{lnx}$.
 
 Étudier très précisément les variations de g (dérivées, sens de variation, concavité, limites, 
asymptotes éventuelles, tableau de variation, graphe). 

Pour toute la suite du problème, on admettra l'existence d'une primitive $G$ de $g$, qu'on ne 
cherchera pas à calculer explicitement. 
 
2) On considère l'ensemble $D = ]0,1/2[\cup ]1,+\infty[$. 
On définit l'intégrale J(x) par : 
$J(x) =\int_{x}^{2x}g(t)dt$

Montrer que $J(x)$ existe pour tout $x\in D$. 

3) Soit l'ensemble $D^+ = [0,1/2[\cup ]1,+\infty[ = D \cup \lbrace 0\rbrace$

On définit l'application $f : D^+\longrightarrow \mathbb{R}$ par : 

$\forall x \in D,f(x) = J(x)$

$f(0)=0$

Montrer que $f$ est dérivable sur $D$, et calculer sa dérivée $f^{\prime}$.

Étudier le signe de $f^{\prime}$ et en déduire le sens des variations de $f$ . 

4) Démontrer que, $\forall x \in D$ : 

$\frac{x}{ln(2x)}\leq f(x)\leq \frac{x}{lnx}$

En déduire les limites de $f(x)$ et de $f(x) /x$ quand $x \longrightarrow 0$. 

Étudier la continuité et la dérivabilité de $f$ en $0$. 

5) Quelles sont les limites de $f(x)$ et de $f(x)/x$ quand $x\longrightarrow \infty$? 

6) Soit l'application $h : ]0,1] \longrightarrow \mathbb{R}$, définie par : 
$h(u) = ln(u)-2u + 2$ 

Montrer qu'il existe un réel unique, noté $\alpha, 0 < \alpha  <1/2$, tel que $h(\alpha) = 0$. 

Montrer que $\forall u \in [\alpha, 1]$, on a $\ln(u)\geq 2u-2$. 

En déduire que $f(x)$ est majorée, pour tout $x\in[\alpha,1/2]$, par $\frac{1}{2}\ln\left(\frac{2x-1}{x-1}\right)$

Calculer la limite de $f$ quand  $x \longrightarrow 1/2$. 

7) Montrer que, $\forall u\geq 1, ln(u) \leq u - 1$. 

En déduire la limite de $f$ quand  $x \longrightarrow 1$. 

8) Construire le tableau de variations de $f$ . 

\medskip\textbf{Exercice 3}

Les symboles $Ln$ et $tan$ représentent respectivement le logarithme népérien et la tangente. 

On donne $Ln2 = 0,69$. 

1) Comparer les intégrales $A$ et $B$ :

$A=\int_{0}^{\frac{\pi}{4}}Ln(\cos x)dx$

$B= \int_{0}^{\frac{\pi}{4}}Ln(\cos(\frac{\pi}{4}- x)dx$.

2. Calculer l'intégrale $I=\int_{0}^{\frac{\pi}{4}}Ln(1+\tan x)dx$.

\end{spacing}
\end{document}
